%% =====================================================================
%%  Chapitre 3 — Variables
%%  Tout le code est vérifié par outils/verifier_code.py (Shell IDLE simulé).
%% =====================================================================
\chapitre{3}{Variables}{Give names to your values}{illusVariables}

\histoire{In chapter~2, every number had to be typed straight into the calculations. Change one
price, and you had to hunt for it everywhere in your program. \textbf{Variables} solve this problem:
they let you give a name to a value, then use the name instead of the value. They are the memory of
your programs. In this chapter, you learn to create and update variables, to choose good names, and
to picture what really happens inside the computer when you write \py{x = 5}.}

\begin{objectifs}
\item create a variable and use it in calculations;
\item read \py{=} correctly: \enquote{the name on the left now refers to the value on the right};
\item update a variable, including with \py{+=}, \py{-=} and \py{*=};
\item choose names that are valid and clear;
\item picture variables as names attached to values in memory;
\item follow the values of variables line by line with a trace table;
\item swap two variables in a single line.
\end{objectifs}

\auprogramme
  {variable, assignment, names and objects, trace table}
  {\py{=}, \py{+=}, \py{-=}, \py{*=}, \py{type()}, multiple assignment}
  {trace tables and memory diagrams}
  {meaningful names in \py{snake\_case}; constants in capitals}
  {a holiday planner you can update in one line}

\seances[4]{livre/otNavy/Name your values/Why variables and how assignment works (3.1--3.2),
  engrenage/otViolet/Inside the memory/Names and objects and the rules for names (3.3--3.4),
  boucle/otTeal/Change and check/Updating variables and types and errors (3.5--3.8),
  crayon/otGreenDark/Practise/Mini-project and exercises and review (3.9--3.11)}

%% ---------------------------------------------------------------------
\rubrique{learn}{Why variables?}

Remember the holiday budget of chapter~2: the number of travellers, \py{2}, appeared seven times in
the program. A \textbf{variable} lets you write it only once, under a name:

\begin{console}
>>> travellers = 2
>>> travellers
2
>>> travellers * 89.5
179.0
\end{console}

The first line creates a variable named \py{travellers} and gives it the value \py{2}. Notice that
Python displays nothing: it has simply memorised the value. From then on, each time you write the
name \py{travellers}, Python replaces it with its value. The instruction \py{travellers = 2} is
called an \textbf{assignment}.

\begin{definition}{variable}
A name that refers to a value stored in the computer's memory. You create a variable by
\textbf{assigning} it a value with the sign \py{=}: \py{price = 4.5}.
\end{definition}

\soustitre{Variables in a program}
In a program, variables make calculations readable: each value has a name that says what it is.

\begin{code}{file=cinema.py}
# cinema.py - the price of an evening at the cinema
ticket_price = 9.5
popcorn_price = 4
people = 3
total = people * (ticket_price + popcorn_price)
print("Total for the evening:", total)
\end{code}
\begin{sortie}
Total for the evening: 40.5
\end{sortie}

Read line~5 aloud: \enquote{the total is the number of people times the price of a ticket plus a
popcorn}. The calculation explains itself. And if the price of a ticket changes, you only change line~2.

\begin{regle}[The assignment]
\py{name = expression}\quad Python first \textbf{evaluates the expression on the right}, then
\textbf{attaches the name on the left} to the result. The name is always on the left, alone.
\end{regle}

%% ---------------------------------------------------------------------
\rubrique{learn}{Assignment step by step}

\soustitre{The right-hand side first}
In \py{total = people * (ticket\_price + popcorn\_price)}, Python does nothing with \py{total} until
it has calculated the right-hand side. The right-hand side can use any variables that already exist.

\soustitre{Changing the value of a variable}
A variable can change its value as often as you like: that is why it is called a \emph{variable}. A
new assignment replaces the old value, which is forgotten:

\begin{console}
>>> score = 10
>>> score = 15
>>> score
15
\end{console}

\soustitre{\texorpdfstring{\py{x = x + 1}}{x = x + 1} is not an equation}
In mathematics, $x = x + 1$ is impossible. In Python, it is one of the most common instructions! Read
it as an assignment: \enquote{calculate \py{x + 1} with the current value of \py{x}, then give the
result to \py{x}.}

\begin{console}
>>> score = 10
>>> score = score + 5
>>> score
15
\end{console}

\begin{attention}[The sign = does not mean \enquote{equals}]
In Python, \py{=} means \enquote{give this value to this name}. That is why the name must be on the
left: \py{10 = x} makes no sense. Python even guesses what you may have meant (the sign \py{==}, which
compares two values, is for chapter~7).
\begin{console}
>>> 10 = x
SyntaxError: cannot assign to literal here. Maybe you meant '==' instead of '='?
\end{console}
\end{attention}

\soustitre{Following a program with a trace table}
To understand a program with several variables, programmers draw a \textbf{trace table}: one line
per instruction, one column per variable, and the value of each variable after the instruction.
Here is a short program and its trace table; the value that changes at each step is highlighted.

\begin{gauche}[0.4]
\begin{code}{file=trace.py}
a = 5
b = a * 2
a = a + 1
b = b - a
print(a, b)
\end{code}
\end{gauche}\begin{droite}[0.6]
\begin{suivi}{2}
\otTH{Line} & \otTH{Instruction} & \otTH{a} & \otTH{b}\\\midrule
1 & \py{a = 5} & \nv{5} & \\
2 & \py{b = a * 2} & \av{5} & \nv{10}\\
3 & \py{a = a + 1} & \nv{6} & \av{10}\\
4 & \py{b = b - a} & \av{6} & \nv{4}\\
\end{suivi}
\end{droite}
\rattacher
\begin{sortie}
6 4
\end{sortie}

At line~2, \py{b} receives \py{10}, calculated with the value of \py{a} at that moment. When \py{a}
changes at line~3, \py{b} does not change: it remembers \py{10}, not the calculation \py{a * 2}.

\begin{atoi}[Your own trace table]
Draw the trace table of this program and predict what it displays. Then check with Python.
\begin{code}{file=trace2.py}
x = 3
y = x + 4
x = y * 2
y = x - y
print(x, y)
\end{code}
\end{atoi}

%% ---------------------------------------------------------------------
\rubrique{hood}{Names and objects}

What really happens in the computer's memory when you write \py{age = 15}? Python creates an
\textbf{object} that represents the integer 15, then attaches the \textbf{name} \py{age} to it, like
a label tied to a box with a string.

\begin{memoire}[after \py{age = 15}]
\nomvar{age}{(0,0)}{n}
\objet{int}{15}{(3.6,0)}{o}
\lien{n}{o}
\end{memoire}

Every value in Python is an object, and every object has a \textbf{type} (here \py{int}) and a
\textbf{value} (here 15). A variable is not a box that contains the value: it is a \textbf{name} that
\emph{refers to} an object.

\soustitre{A new assignment moves the label}
When you write \py{age = 16}, Python creates a new object, 16, and moves the label \py{age} onto it.
The object 15 is no longer attached to any name: Python will automatically free the memory it uses.

\begin{memoire}[after \py{age = 16}]
\nomvar{age}{(0,0)}{n}
\objet{int}{16}{(3.6,0)}{o}
\objet{int}{15}{(3.6,1.6)}{a}
\lien{n}{o}
\draw[otLienAncien] (n.north east) -- (a.west);
\node[anchor=west,font=\otRegular\fontsize{8}{9}\selectfont,text=otMuted] at (a.east) {\ forgotten};
\end{memoire}

\soustitre{Two names for the same object}
An assignment can attach a second name to an object that already exists. After \py{a = 42} then
\py{b = a}, the two names refer to the \emph{same} object:

\begin{memoire}[after \py{a = 42} and \py{b = a}]
\nomvar{a}{(0,0.65)}{na}
\nomvar{b}{(0,-0.65)}{nb}
\objet{int}{42}{(3.6,0)}{o}
\lien{na}{o}
\lien{nb}{o}
\end{memoire}

What happens to \py{b} if you now give a new value to \py{a}? The label \py{a} moves to a new object,
and \py{b} stays where it was:

\begin{console}
>>> a = 42
>>> b = a
>>> a = 43
>>> b
42
\end{console}

\begin{memoire}[after \py{a = 43}]
\nomvar{a}{(0,0.65)}{na}
\nomvar{b}{(0,-0.65)}{nb}
\objet{int}{43}{(3.6,0.65)}{o1}
\objet{int}{42}{(3.6,-0.65)}{o2}
\lien{na}{o1}
\lien{nb}{o2}
\end{memoire}

\begin{regle}[Names and objects]
A variable is a name attached to an object. An assignment never copies or modifies an object: it
attaches a name to it. Reassigning a name moves the label, and the other names are not affected.
\end{regle}

\begin{culture}[The type belongs to the object]
In Python, a variable has no fixed type: the same name can refer to a number, then to a piece of
text. It is the \emph{object} that has a type. This is called \textbf{dynamic typing}. It makes
Python flexible, but it also means you must choose your names carefully, so that you always know what
they refer to.
\end{culture}

%% ---------------------------------------------------------------------
\rubrique{learn}{Naming variables}

\soustitre{The rules}
Python accepts a name only if it respects three rules:
\begin{itemize}
\item it contains only letters, digits and underscores \py{\_}; no spaces, no hyphens, no other
  signs;
\item it does not start with a digit: \py{place2} is valid, \py{2nd\_place} is not;
\item it is not one of Python's \textbf{keywords}, the words reserved by the language.
\end{itemize}
Python is case-sensitive: \py{age}, \py{Age} and \py{AGE} are three different names.

\begin{console}
>>> 2nd_place = "Ana"
SyntaxError: invalid decimal literal
>>> my name = "Ana"
SyntaxError: invalid syntax
>>> user-name = "Ana"
SyntaxError: cannot assign to expression here. Maybe you meant '==' instead of '='?
>>> class = 3
SyntaxError: invalid syntax
\end{console}

The third message is surprising: Python reads \py{user-name} as the calculation \py{user minus name},
and you cannot assign a value to a calculation. The last one fails because \py{class} is a keyword.

\begin{encadre}{cle}{Python's 35 keywords}[You cannot use them as names]
\begin{tabularx}{\linewidth}{@{}*{7}{>{\ttfamily\color{otCodeEncre}}X}@{}}
False & None & True & and & as & assert & async\\
await & break & class & continue & def & del & elif\\
else & except & finally & for & from & global & if\\
import & in & is & lambda & nonlocal & not & or\\
pass & raise & return & try & while & with & yield\\
\end{tabularx}
\end{encadre}

You don't need to learn this list: IDLE colours keywords in orange as soon as you type them, which
warns you immediately.

\soustitre{The conventions}
A valid name is not necessarily a good name. Python programmers follow the conventions of PEP~8:
\begin{itemize}
\item names of variables in lower case, words separated by underscores: \py{ticket\_price},
  \py{number\_of\_guests}. This style is called \textbf{snake\_case};
\item names that say what the value represents: \py{price\_per\_night} rather than \py{p} or
  \py{x2};
\item for a value that should never change in the program, a \textbf{constant}, capital letters:
  \py{VAT\_RATE = 0.2}, \py{SECONDS\_PER\_DAY = 86400}.
\end{itemize}

\begin{tcolorbox}[otBase,colback=white,colframe=otRule,boxrule=0.5pt,left=4mm,right=4mm]
\renewcommand{\arraystretch}{1.3}%
\begin{tabularx}{\linewidth}{@{}p{34mm}p{22mm}Y@{}}
\otTH{Name} & \otTH{Valid?} & \otTH{Comment}\\\midrule
\py{price\_per\_night} & \otCheck & perfect: clear, in snake\_case\\
\py{pricePerNight} & \otCheck & valid, but not the Python style (this is \enquote{camelCase})\\
\py{p} & \otCheck & valid, but says nothing\\
\py{price per night} & \otCroix & spaces are forbidden\\
\py{price-per-night} & \otCroix & hyphens are read as minus signs\\
\py{3rd\_night} & \otCroix & starts with a digit\\
\py{MAX\_GUESTS} & \otCheck & a constant\\
\end{tabularx}
\end{tcolorbox}

%% ---------------------------------------------------------------------
\rubrique{learn}{Updating variables}

Programs constantly update variables: a score increases, a balance decreases, a price goes up by a
percentage. Python offers a shortcut for these updates: \py{score += 5} means exactly the same as
\py{score = score + 5}.

\begin{console}
>>> balance = 100
>>> balance -= 30
>>> balance += 50
>>> balance *= 1.02
>>> balance
122.4
\end{console}

\begin{suivi}[What happened to the balance]{1}
\otTH{Line} & \otTH{Instruction} & \otTH{balance}\\\midrule
1 & \py{balance = 100} & \nv{100}\\
2 & \py{balance -= 30} \hspace{3mm}{\color{otMuted}withdraw 30} & \nv{70}\\
3 & \py{balance += 50} \hspace{3mm}{\color{otMuted}deposit 50} & \nv{120}\\
4 & \py{balance *= 1.02} \hspace{3mm}{\color{otMuted}add 2\,\% interest} & \nv{122.4}\\
\end{suivi}

\begin{tcolorbox}[otBase,colback=white,colframe=otRule,boxrule=0.5pt,left=4mm,right=4mm]
\otEtiquette{otNavy}{Augmented assignment}\par\vspace{1.4mm}
\renewcommand{\arraystretch}{1.25}%
\begin{tabularx}{\linewidth}{@{}p{30mm}p{40mm}Y@{}}
\py{x += 3} & \py{x = x + 3} & add\\
\py{x -= 3} & \py{x = x - 3} & subtract\\
\py{x *= 3} & \py{x = x * 3} & multiply\\
\py{x /= 3} & \py{x = x / 3} & divide (gives a float)\\
\py{x //= 3}\quad\py{x \%= 3} & \py{x = x // 3}\quad\py{x = x \% 3} & quotient, remainder\\
\py{x **= 3} & \py{x = x ** 3} & power\\
\end{tabularx}
\end{tcolorbox}

\begin{definition}{augmented assignment}
An instruction such as \py{x += 3} that updates a variable using its current value. It is a shorter
way of writing \py{x = x + 3}.
\end{definition}

\begin{atoi}[A game of points]
A player starts with 0 points. She wins 50 points, loses 20, then her score is tripled by a bonus.
Write the program with \py{+=}, \py{-=} and \py{*=}, and display the final score.
\end{atoi}

%% ---------------------------------------------------------------------
\rubrique{learn}{Variables and types}

A variable can refer to any kind of value: an integer, a float, but also a piece of text, which
Python calls a \py{str} (for \emph{string}), or a truth value, \py{True} or \py{False}, of type
\py{bool}. You will study strings in chapter~4 and truth values in chapter~7.

\begin{code}{file=profile.py}
name = "Ana"           # str: a piece of text
age = 15               # int: a whole number
height = 1.62          # float: a decimal number
likes_python = True    # bool: True or False
print(name, "is", age, "years old and", height, "m tall.")
print("Next year,", name, "will be", age + 1)
print(type(name), type(age), type(height), type(likes_python))
\end{code}
\begin{sortie}
Ana is 15 years old and 1.62 m tall.
Next year, Ana will be 16
<class 'str'> <class 'int'> <class 'float'> <class 'bool'>
\end{sortie}

Notice the difference between \py{"Ana"}, with quotes, which is a piece of text, and \py{name},
without quotes, which is the name of a variable. \py{print(name)} displays \py{Ana};
\py{print("name")} displays the word \py{name}.

\begin{attention}[One name, one meaning]
Python lets you write \py{x = 5} and then \py{x = "five"}: the same name now refers to a string. Your
program becomes very hard to follow. Give each variable one role, and keep the same type for it from
beginning to end.
\end{attention}

%% ---------------------------------------------------------------------
\rubrique{debug}{Errors with variables}

\cas{1}{A typing mistake in a name}
Python cannot guess that \py{prise} is a typo for \py{price}: for it, it is another name, which does
not exist. Fortunately, it suggests the closest existing name.
\begin{console}
>>> price = 4.5
>>> prise * 2
Traceback (most recent call last):
  File "<pyshell#1>", line 1, in <module>
    prise * 2
NameError: name 'prise' is not defined. Did you mean: 'price'?
\end{console}

\cas{2}{The wrong case}
\begin{console}
>>> price = 4.5
>>> Price * 2
Traceback (most recent call last):
  File "<pyshell#1>", line 1, in <module>
    Price * 2
NameError: name 'Price' is not defined. Did you mean: 'price'?
\end{console}

\cas{3}{A variable used before it exists}
Python runs a program from top to bottom. A variable only exists from the line where it is assigned:
using it earlier causes a \py{NameError}.
\begin{code}{file=too_early.py, bug}
print("Total:", total)
total = 3 * 4.5
\end{code}
\begin{erreur}
Traceback (most recent call last):
  File "too_early.py", line 1, in <module>
    print("Total:", total)
NameError: name 'total' is not defined
\end{erreur}

\cas{4}{A variable that does not follow}
This is not an error message but a misunderstanding, very common among beginners. The instruction
\py{total = price * 3} gives \py{total} a number, calculated once and for all: if \py{price} changes
afterwards, \py{total} does not change.
\begin{code}{file=not_updated.py}
price = 4
total = price * 3
price = 5
print("Total:", total)
\end{code}
\begin{sortie}
Total: 12
\end{sortie}
To update the total, you must calculate it again after changing the price.

%% ---------------------------------------------------------------------
\rubrique{pro}{Good names and handy tricks}

\soustitre{Choose names for the reader}
Compare these two programs. They do exactly the same thing:

\begin{gauche}
\begin{code}{file=bad_names.py}
a = 12.5
b = 35
c = a * b
print(c)
\end{code}
\end{gauche}\begin{droite}
\begin{code}{file=good_names.py}
hourly_rate = 12.5
hours = 35
weekly_pay = hourly_rate * hours
print(weekly_pay)
\end{code}
\end{droite}

The second one needs no comment: the names explain everything. Good names are the best
documentation. Avoid single letters (except in very short calculations), meaningless names such as
\py{data} or \py{thing}, and the letters \py{l}, \py{O} and \py{I} on their own, which look like the
digits \py{1} and \py{0}.

\soustitre{Several assignments in one line}
Python can assign several variables at once, separating the names and the values with commas:

\begin{console}
>>> width, height = 8, 5
>>> width * height
40
\end{console}

\soustitre{Swapping two variables}
To swap the values of two variables, the classic method uses a third, temporary variable: you cannot
write \py{a = b} then \py{b = a}, because the first assignment has already lost the value of
\py{a}. Python offers a much more elegant method:

\begin{console}
>>> a, b = 1, 2
>>> a, b = b, a
>>> a, b
(2, 1)
\end{console}

Python first evaluates the whole right-hand side, \py{b, a}, that is \py{2, 1}; only then does it
assign the two values to \py{a} and \py{b}. (The parentheses in the answer of the Shell indicate a
\emph{tuple}, a group of values that you will study in chapter~19.)

\begin{pratique}[Constants at the top]
Put the values your program depends on at the top of the file, with clear names, and constants in
capitals: \py{VAT\_RATE = 0.2}. Anyone who reads the program sees at once what can be adjusted, and
a change happens in one single place.
\end{pratique}

%% ---------------------------------------------------------------------
\rubrique{build}{Mini-project: a smarter holiday planner}

Let's rewrite the holiday budget of chapter~2 with variables. All the data go at the top of the
program, under clear names; the calculations use only names; the display comes last. This
three-part structure\,---\,\textbf{data, calculations, display}\,---\,is that of countless
professional programs.

\begin{code}{file=holiday_planner.py}
# holiday_planner.py - the Lisbon budget, now with variables

# Data
travellers = 2
nights = 6
days = 7
flight_price = 89.5      # return ticket, per person
room_price = 74          # per night, for the room
food_per_day = 25        # per person
museum_pass = 21.5       # per person
months = 6               # months left to save

# Calculations
flights = travellers * flight_price
hotel = nights * room_price
food = travellers * days * food_per_day
museums = travellers * museum_pass
total = flights + hotel + food + museums

# Display
print("=== HOLIDAY IN LISBON ===")
print("Travellers:", travellers)
print("Flights:", flights)
print("Hotel:", hotel)
print("Food:", food)
print("Museums:", museums)
print("Total:", total)
print("Per person:", round(total / travellers, 2))
print("To save each month:", round(total / months, 2))
\end{code}
\begin{sortie}
=== HOLIDAY IN LISBON ===
Travellers: 2
Flights: 179.0
Hotel: 444
Food: 350
Museums: 43.0
Total: 1016.0
Per person: 508.0
To save each month: 169.33
\end{sortie}

The results are exactly those of chapter~2. But now, if Leo's sister joins the trip, you change
\textbf{one single line}, and every result follows:

\begin{code}{file=holiday_planner.py, variante, first=4, marks={1}}
travellers = 3
\end{code}
\begin{sortie}
=== HOLIDAY IN LISBON ===
Travellers: 3
Flights: 268.5
Hotel: 444
Food: 525
Museums: 64.5
Total: 1302.0
Per person: 434.0
To save each month: 217.0
\end{sortie}

The hotel does not change, since the room price is per night: the program reflects the real rules
of the budget, written once and for all.

\begin{atoi}[One more night]
The friends decide to stay one more night. Which variables must you change? Make the change and run
the program: what is the new total for three travellers?
\end{atoi}

%% ---------------------------------------------------------------------
\rubrique{practice}{Exercises}

\exercice[type=read,niveau=1]{Fill in the trace table}{Copy and complete the trace table of this
program, then write what it displays.}
\begin{code}{file=trace3.py}
a = 4
b = a * 3
a = b - a
b = b // a
print(a, b)
\end{code}
\begin{suivi}{2}
\otTH{Line} & \otTH{Instruction} & \otTH{a} & \otTH{b}\\\midrule
1 & \py{a = 4} & \trou[12mm]{} & \trou[12mm]{}\\
2 & \py{b = a * 3} & \trou[12mm]{} & \trou[12mm]{}\\
3 & \py{a = b - a} & \trou[12mm]{} & \trou[12mm]{}\\
4 & \py{b = b // a} & \trou[12mm]{} & \trou[12mm]{}\\
\end{suivi}

\exercice[type=predict,niveau=1]{Does it follow?}{What does this program display? Explain why
\py{y} does not change at line~3.}
\begin{code}{file=follow.py}
x = 10
y = x
x = x + 5
print(x, y)
\end{code}
\ecranvide{1}

\exercice[type=explain,niveau=1]{Valid or not?}{For each name, say whether Python accepts it. If it
does not, explain why; if it does, say whether it follows the Python conventions.}
\begin{questions}[2]
\item \py{age}
\item \py{2fast}
\item \py{first\_name}
\item \py{first-name}
\item \py{\_count}
\item \py{class}
\item \py{Class}
\item \py{my score}
\item \py{totalPrice}
\item \py{MAX\_SPEED}
\end{questions}

\exercice[type=write,niveau=1]{The rectangle}{A rectangle is 7.5~m wide and 4~m high. Write a
program that stores these two measurements in variables, calculates the area and the perimeter, and
displays them as below.}
\cible{rectangle}
\begin{sortie}
Area: 30.0 square metres
Perimeter: 23.0 metres
\end{sortie}

\exercice[type=complete,niveau=2]{The shopping list}{Complete the program so that it displays the
output below. A loaf of bread costs 1.25 euros and a bottle of milk 0.75 euros.}
\begin{code}{file=shopping.py, trous}
bread_price = 1.25
milk_price = ________
milk_bottles = 3
total = bread_price + __________ * milk_bottles
print("Total:", _____)
\end{code}
\cible{courses}
\begin{sortie}
Total: 3.5
\end{sortie}

\exercice[type=fix,niveau=2]{The weekly pay}{This program should calculate a weekly pay, but it
stops with an error. Correct it: there are two mistakes.}
\begin{code}{file=weekly_pay.py, bug}
hourly_rate = 12.5
hours = 35
weekly_pay = hourly_rate * hour
print("Weekly pay:", Weekly_pay)
\end{code}
\begin{erreur}
Traceback (most recent call last):
  File "weekly_pay.py", line 3, in <module>
    weekly_pay = hourly_rate * hour
NameError: name 'hour' is not defined. Did you mean: 'hours'?
\end{erreur}
\cible{salaire}
\begin{sortie}
Weekly pay: 437.5
\end{sortie}

\exercice[type=write,niveau=2]{Swap!}{The variable \py{left} contains \py{"red"} and the variable
\py{right} contains \py{"blue"}. Write a program that swaps their values, first with a temporary
variable, then with Python's one-line method, and displays the result as below.}
\cible{echange}
\begin{sortie}
Before: red blue
After the first swap: blue red
After the second swap: red blue
\end{sortie}

\exercice[type=write,niveau=2]{Average speed}{A train travels 150~km in 2.5~hours. Write a program
that stores the distance and the duration in variables and displays the average speed, in km/h and
in metres per second, rounded to the hundredth. (To convert km/h into m/s, divide by 3.6.)}
\cible{vitesse}
\begin{sortie}
Average speed: 60.0 km/h
That is 16.67 m/s
\end{sortie}

\exercice[type=predict,niveau=2]{Augmented assignments}{Draw the trace table of this program and
predict what it displays.}
\begin{code}{file=points.py}
points = 10
points += 5
points *= 2
points -= 4
points //= 3
print(points)
\end{code}
\ecranvide{1}

\exercice[type=write,niveau=2]{Pancakes for six}{A pancake recipe for 4 people needs 250~g of flour,
500~ml of milk and 2 eggs. Write a program that calculates the quantities for any number of guests:
store the number of guests in a variable, calculate a multiplying factor, and display the quantities
for 6 guests as below. Then check that it works for 2 guests.}
\cible{crepes}
\begin{sortie}
Pancakes for 6 people:
Flour: 375.0 g
Milk: 750.0 ml
Eggs: 3.0
\end{sortie}

\exercice[type=read,niveau=2]{Draw the memory}{Draw the memory diagram (names and objects) after
each line of this program. Then say what \py{print(a, b)} would display.}
\begin{code}{file=memory.py}
a = 3
b = a
a = 7
\end{code}

\exercice[type=challenge,niveau=3]{Digit games}{The variable \py{n} contains the number 4728. Using
\py{//} and \py{\%}, store each of its four digits in a variable (\py{thousands}, \py{hundreds},
\py{tens}, \py{units}), then display the sum of its digits and the number written backwards. Your
program must still work if you replace 4728 by any other four-digit number.}
\cible{chiffres}
\begin{sortie}
Sum of the digits: 21
Backwards: 8274
\end{sortie}

%% ---------------------------------------------------------------------
\rubrique{review}{Chapter review}

\begin{essentiel}
\otEssTitre{Creating and using a variable}
\py{name = expression}: Python evaluates the right-hand side, then attaches the name to the result.
The name can then be used in any expression.
\begin{code}{mini}
price = 4.5
quantity = 3
total = price * quantity   # 13.5
\end{code}

\otEssTitre{Changing a value}
A new assignment replaces the value; the old one is forgotten. \py{x = x + 1}, or \py{x += 1}, means
\enquote{the new value of \py{x} is its old value plus 1}. Augmented assignments: \py{+=}, \py{-=},
\py{*=}, \py{/=}, \py{//=}, \py{\%=}, \py{**=}.

\otEssTitre{Names and objects}
A variable is a \textbf{name attached to an object}; the object has a type and a value. Several names
can refer to the same object; reassigning one name never changes the others. A variable computed from
others does not follow their later changes.

\otEssTitre{Names}
Letters, digits and \py{\_}; never a digit first; no keyword; case matters. Convention:
\py{snake\_case} for variables, \py{UPPER\_CASE} for constants, and always meaningful names.

\otEssTitre{Two handy tricks}
\py{width, height = 8, 5} assigns two variables at once; \py{a, b = b, a} swaps them.
\end{essentiel}

\begin{quiz}
\qcm{After \py{x = 4} and \py{x = x * 3}, what is the value of \py{x}?}{\py{4}}{\py{12}}{\py{7}}
\qcm{What does \py{price = 9.5} display in the Shell?}{nothing}{\py{9.5}}{\py{price = 9.5}}
\qcm{Which name is valid?}{\py{2nd\_try}}{\py{second-try}}{\py{second\_try}}
\qcm{What is the value of \py{n} after \py{n = 10}, then \py{n -= 3}?}{\py{3}}{\py{7}}{\py{-3}}
\begin{qcmcode}{What does this program display?}{\py{5 5}}{\py{8 5}}{\py{8 8}}
\begin{code}{mini}
a = 5
b = a
a = a + 3
print(a, b)
\end{code}
\end{qcmcode}
\qcm{Which instruction swaps \py{a} and \py{b}?}{\py{a = b}}{\py{a, b = b, a}}{\py{a == b}}
\qcm{What is the convention for the name of a constant?}{\py{vat\_rate}}{\py{VatRate}}{\py{VAT\_RATE}}
\qcm{What is the type of \py{"42"}?}{\py{int}}{\py{str}}{\py{float}}
\qcm{A program uses \py{total} on line~2 and creates it on line~5. What happens?}{it works}{a
  \py{NameError}}{a \py{SyntaxError}}
\qcm{Which name is the most readable?}{\py{x}}{\py{np}}{\py{number\_of\_players}}
\end{quiz}

\begin{jesais}
\item create a variable and use it in an expression;
\item explain what \py{x = x + 1} does;
\item update a variable with \py{+=}, \py{-=}, \py{*=} and the other augmented assignments;
\item draw a trace table and a memory diagram for a short program;
\item choose valid names that follow the conventions;
\item recognise and fix a \py{NameError};
\item assign several variables at once and swap two variables.
\end{jesais}
